%%%%%%%%%%%%%%%%%
% CV tailored for Thales - Ingénieur Intégration Vérification Validation Logiciel
%%%%%%%%%%%%%%%%%

\documentclass[8pt,a4paper,ragged2e,withhyper]{altacv}

\geometry{left=1.25cm,right=1.25cm,top=.5cm,bottom=1.cm,columnsep=1.2cm}

\usepackage{paracol}
\usepackage{fontawesome5}
\usepackage{hyperref}

\iftutex 
  \setmainfont{Roboto Slab}
  \setsansfont{Lato}
  \renewcommand{\familydefault}{\sfdefault}
\else
  \usepackage[rm]{roboto}
  \usepackage[defaultsans]{lato}
  \renewcommand{\familydefault}{\sfdefault}
\fi

\definecolor{SlateGrey}{HTML}{2E2E2E}
\definecolor{LightGrey}{HTML}{666666}
\definecolor{DarkPastelRed}{HTML}{8AC0D9}
\definecolor{PastelRed}{HTML}{00B0DA}
\definecolor{GoldenEarth}{HTML}{B4AD0C}

\colorlet{name}{black}
\colorlet{tagline}{PastelRed}
\colorlet{heading}{DarkPastelRed}
\colorlet{headingrule}{GoldenEarth}
\colorlet{subheading}{PastelRed}
\colorlet{accent}{PastelRed}
\colorlet{emphasis}{SlateGrey}
\colorlet{body}{LightGrey}

\definecolor{violet}{HTML}{5A2582}
\definecolor{rouge}{HTML}{F53B3B}
\definecolor{noir}{HTML}{00232C}
\definecolor{bleu}{HTML}{00B0DA}
\definecolor{rose}{HTML}{F831A0}
\definecolor{jaune}{HTML}{FFEC00}
\definecolor{vert}{HTML}{34AD6C}

\renewcommand{\namefont}{\Huge\rmfamily\bfseries}
\renewcommand{\personalinfofont}{\footnotesize}
\renewcommand{\cvsectionfont}{\LARGE\rmfamily\bfseries}
\renewcommand{\cvsubsectionfont}{\large\bfseries}
\renewcommand{\cvItemMarker}{{\small\textbullet}}
\renewcommand{\cvRatingMarker}{\faCircle}

\input{pubs-num.tex}
\addbibresource{sample.bib}

\begin{document}

\name{BOILEAU Guillaume}
\tagline{Ingénieur IVVQ Logiciel | Intégration, validation, simulation, Linux \& traitement du signal}

\photoL{2.8cm}{moi}

\personalinfo{%
  \email{guillaume.boileaupro@gmail.com}
  \phone{+33 6 59 94 85 84}
  \location{Alpes-Maritimes, FRANCE}
  \linkedin{guillaume-boileau-891b81265}
  \researchgate{Guillaume-Boileau-2}
  \orcid{0000-0002-3576-6968}
}

\makecvheader
\vspace{-0.3cm}

\AtBeginEnvironment{itemize}{\small}
\columnratio{0.64}

\begin{paracol}{2}

\cvsection{Experience}

\cvevent{\textbf{Ingénieur logiciel full stack -- Intégration \& support opérationnel}}{VEV Platform Services France}{2025 -- 2026}{Valbonne, France}

\textbf{Contexte :} Développement, intégration et support opérationnel pour une plateforme CPMS liée à des infrastructures de recharge de véhicules électriques.

\textbf{Objectif :} Contribuer au développement, à l’intégration et à la fiabilisation de services logiciels connectés à des équipements terrain, des flux opérationnels et des interfaces système.

\begin{itemize}
  \item \textbf{Développement logiciel :} Contribution au développement d’applications web et de services backend en TypeScript, Angular et Node.js.
  \item \textbf{Intégration applicative :} Intégration de services logiciels, données opérationnelles, interfaces applicatives et équipements connectés.
  \item \textbf{Vérification fonctionnelle :} Contrôle de la cohérence entre données applicatives, comportement des équipements et attentes opérationnelles.
  \item \textbf{Analyse d’anomalies :} Investigation de dysfonctionnements applicatifs, analyse de comportements inattendus et contribution à l’identification des causes racines.
  \item \textbf{Support production :} Suivi technique des incidents, analyse de logs et contribution à la fiabilisation des services en exploitation.
  \item \textbf{Flux opérationnels :} Analyse de flux de données, notamment via des échanges FTP utilisés pour le pilotage et le suivi technique.
  \item \textbf{Monitoring \& outillage :} Utilisation d’outils de suivi, logs et diagnostic tels que OpenSearch, Sentry, GitLab, Confluence et Zenhub.
  \item \textbf{Coordination agile :} Suivi des tâches, priorisation, backlog, documentation et coordination dans un environnement itératif.
\end{itemize}

\divider

\cvevent{\textbf{Ingénieur R\&D senior -- IVVQ, simulation \& validation logicielle}}{Laboratoire Lagrange, Observatoire de la Côte d’Azur, CNES}{2023 -- 2025}{Nice, France}

\textbf{Contexte :} Activités d’ingénierie et de R\&D pour la mission spatiale LISA ESA/NASA, dans un environnement international impliquant simulation, intégration logicielle, traitement de données et validation technique.

\textbf{Objectif :} Développer et structurer des outils Python permettant d’intégrer, comparer, vérifier, valider et documenter des contributions techniques hétérogènes.

\begin{itemize}
  \item \textbf{IVVQ logiciel :} Mise en place de contrôles de cohérence, règles de comparaison, vérifications de résultats et méthodes de validation technique.
  \item \textbf{Intégration de composants :} Consolidation de modèles, données, sorties de simulation et contributions logicielles dans un cadre commun d’analyse.
  \item \textbf{Validation sur simulation :} Vérification de résultats issus de simulations numériques complexes, comparaison de modèles et analyse de sensibilité.
  \item \textbf{Traitement du signal :} Analyse de signaux faibles, bruit, performances et limites de détection dans un contexte spatial exigeant.
  \item \textbf{Analyse de dysfonctionnements :} Identification d’incohérences, diagnostic de résultats non conformes et contribution à la compréhension des causes techniques.
  \item \textbf{Stratégie de test :} Structuration de workflows de validation, configurations, versions, sorties de simulation et résultats d’analyse.
  \item \textbf{Coordination logiciel/système :} Interface avec des contributeurs scientifiques, logiciels et institutionnels dans un environnement international.
  \item \textbf{Documentation technique :} Utilisation de Confluence pour organiser la documentation, la traçabilité des analyses et le partage de connaissances.
  \item \textbf{Plateforme logicielle :} Développement de CATpyLISA, plateforme Python pour l’intégration, la comparaison, la validation et la revue technique de modèles.
  \textcolor{DarkPastelRed}{\href{https://gitlab.oca.eu/gboileau/lisa_l3_tools}{\bf GitLab:CATpyLISA}}
\end{itemize}

\divider

\cvevent{\textbf{Data Scientist -- Modélisation, bruit \& analyse de performance}}{Universiteit Antwerpen, Belgium}{2021 -- 2023}{Antwerp, Belgium}

\textbf{Contexte :} Analyse de données, modélisation statistique et études de performance pour des instruments de précision dans un environnement technique international.

\textbf{Objectif :} Développer des workflows robustes et reproductibles pour identifier, caractériser et réduire des sources affectant la qualité de mesure.

\begin{itemize}
  \item \textbf{Analyse de performance :} Évaluation de la qualité de mesure, des limites instrumentales et de l’impact de différentes sources de bruit.
  \item \textbf{Traitement du signal :} Identification, caractérisation et comparaison de contributions de bruit dans des mesures complexes.
  \item \textbf{Modélisation statistique :} Développement de méthodes d’analyse et d’inférence bayésienne, notamment par chaînes MCMC.
  \item \textbf{Pipelines Python :} Mise en place de workflows de traitement, filtrage, transformation, analyse statistique et visualisation.
  \item \textbf{Validation de résultats :} Analyse de sensibilité, vérification de cohérence et évaluation des limites de performance.
  \item \textbf{Reporting technique :} Production de résultats traçables, de synthèses techniques et de supports pour des parties prenantes internationales.
\end{itemize}

\divider

\cvevent{\textbf{Ingénieur R\&D junior -- Signal, simulation \& validation}}{ARTEMIS, Observatoire de la Côte d’Azur}{2018 -- 2021}{Nice, France}

\textbf{Contexte :} Activités de R\&D liées à l’analyse de signaux faibles, à la simulation numérique et à la préparation de missions spatiales et d’instruments de détection.

\textbf{Objectif :} Développer des méthodes d’analyse, automatiser des simulations et produire des résultats exploitables dans un environnement scientifique et technique exigeant.

\begin{itemize}
  \item \textbf{Simulation numérique :} Développement d’approches de simulation pour des environnements de signaux complexes et des jeux de données de grande taille.
  \item \textbf{Analyse de données :} Mise en place de méthodes pour analyser, comparer et séparer des signaux faibles et superposés.
  \item \textbf{Traitement du signal :} Études de caractérisation, de décomposition et de comparaison de signaux dans des contextes bruités.
  \item \textbf{Validation technique :} Vérification de cohérence, comparaison de résultats, analyses d’écarts et études de sensibilité.
  \item \textbf{Investigation technique :} Analyse de sources de bruit et diagnostics par corrélations entre capteurs.
  \item \textbf{Documentation :} Formalisation des méthodes, hypothèses, résultats et conclusions pour revue collaborative.
\end{itemize}

\switchcolumn

\cvsection{Profile}

\textbf{Ingénieur docteur orienté IVVQ logiciel, intégration applicative, validation de systèmes complexes, simulation numérique, Linux, automatisation et traitement du signal.}

Mon parcours combine développement logiciel, intégration de composants, validation sur simulation, analyse de performance, investigation d’anomalies, documentation technique et coordination avec des équipes logicielles, scientifiques et systèmes.

J’ai travaillé dans des environnements exigeants : mission spatiale LISA ESA/NASA, CNES, collaborations internationales, instrumentation scientifique et plateforme logicielle industrielle connectée à des équipements terrain.

\medskip

\textbf{Positionnement cible :} Ingénieur Intégration Vérification Validation Logiciel, ingénieur IVVQ système, ingénieur validation logicielle, ingénieur intégration applicative ou ingénieur test logiciel pour systèmes complexes.

\begin{itemize}
  \item Expérience en intégration logicielle, validation de résultats, contrôle de cohérence et analyse d’anomalies.
  \item Solide culture Linux, Git/GitLab, Docker, Bash, Python, CI/CD et environnements logiciels techniques.
  \item Forte compétence différenciante en traitement du signal, simulation, bruit, signaux faibles et analyse de performance.
  \item Capacité à collaborer avec équipes logicielles, systèmes et parties prenantes techniques dans des environnements complexes.
\end{itemize}

\cvsection{Languages}

\cvskill{English}{4}
\divider
\cvskill{Spanish}{2}
\divider

\medskip

\cvsection{Education}

\cvevent{PhD in Physics}{Université Côte d’Azur}{2018 -- 2021}{Nice, France}

\divider

\cvevent{Selective Graduate Program in Fundamental Physics (Magistère)}{Université Paris-Saclay}{2016 -- 2018}{Orsay, France}

\divider

\cvevent{MSc in Astronomy and Astrophysics}{Observatoire de Paris / Université Paris-Saclay}{2017 -- 2018}{Paris, France}

\divider

\cvevent{BSc in Fundamental Physics}{Université Paris-Saclay}{2015 -- 2016}{Orsay, France}

\cvsection{Professional Training}

\cvevent{Supervised Machine Learning: Regression \& Classification}{DeepLearning.AI -- Andrew Ng}{2026}{Coursera}
\divider

\cvevent{Développement assisté par IA}{Prompting, relecture de code, debugging, prompting incrémental et critique}{2026}{Autoformation}

\end{paracol}

\cvsection{Projects}

\cvevent{CATpyLISA -- Plateforme Python d’intégration, comparaison et validation}{Mission spatiale LISA ESA/NASA}{2023 -- 2025}{}

Développement d’une plateforme Python dédiée à l’intégration de modèles, à la structuration de données, à la comparaison de résultats, à la vérification, à la validation technique et à la revue de workflows liés à la mission LISA.

\textbf{Rôle :} développement Python, conception de pipelines, intégration de contributions, structuration de données, logique de validation, contrôles de cohérence, analyse d’écarts, documentation technique et coordination avec plusieurs contributeurs.

\textbf{Compétences mobilisées :} Python, Linux, GitLab, intégration logicielle, validation, traitement du signal, analyse de performance, simulation, documentation technique, environnement spatial international.

\divider

\cvevent{Workflows de simulation, performance et analyse de bruit}{Instrumentation scientifique et préparation de missions spatiales}{2018 -- 2025}{}

Développement de workflows numériques et statistiques pour analyser des signaux faibles, caractériser des contributions de bruit, comparer des résultats, évaluer la robustesse de modèles et produire des résultats traçables.

\textbf{Rôle :} simulation, traitement du signal, analyse de performance, caractérisation de bruit, études de sensibilité, diagnostics techniques, validation de résultats et reporting.

\textbf{Compétences mobilisées :} traitement du signal, simulation numérique, analyse spectrale, Python, C, validation, analyse de sensibilité, diagnostic technique, documentation.

\cvsection{Compétences clés}

\vspace{-.3cm}

\cvsubsection{IVVQ logiciel, intégration \& validation}

{\small
\cvtag{IVVQ logiciel}
\cvtag{Intégration logicielle}
\cvtag{Vérification}
\cvtag{Validation}
\cvtag{Tests fonctionnels}
\cvtag{Tests de non-régression}
\cvtag{Validation sur simulation}
\cvtag{Validation sur plateforme}
\cvtag{Couverture de test}
\cvtag{Stratégie de test}
\cvtag{Cycle en V}
\cvtag{Méthodes agiles}
\cvtag{Contrôles de cohérence}
\cvtag{Analyse d’écarts}
\cvtag{Traçabilité}
}

\divider\smallskip
\vspace{-.3cm}

\cvsubsection{Analyse d’anomalies, support \& diagnostic}

{\small
\cvtag{Analyse d’anomalies}
\cvtag{Root cause analysis}
\cvtag{Diagnostic technique}
\cvtag{Support production}
\cvtag{Analyse de logs}
\cvtag{Investigation incident}
\cvtag{Suivi de corrections}
\cvtag{Fiabilisation}
\cvtag{OpenSearch}
\cvtag{Sentry}
\cvtag{Confluence}
\cvtag{Jira}
\cvtag{Zenhub}
\cvtag{Reporting technique}
}

\divider\smallskip
\vspace{-.3cm}

\cvsubsection{Logiciel, Linux \& automatisation}

{\small
\cvtag{Python}
\cvtag{C}
\cvtag{Bash}
\cvtag{Shell scripting}
\cvtag{Linux}
\cvtag{Git}
\cvtag{GitLab}
\cvtag{Docker}
\cvtag{CI/CD}
\cvtag{SQL}
\cvtag{TypeScript}
\cvtag{Angular}
\cvtag{Node.js}
\cvtag{REST API}
\cvtag{FTP}
\cvtag{Pandas}
\cvtag{NumPy}
\cvtag{SciPy}
\cvtag{Power BI}
}

\divider\smallskip
\vspace{-.3cm}

\cvsubsection{Outils industriels \& adaptation}

{\small
\cvtag{GitLab}
\cvtag{Git}
\cvtag{Docker}
\cvtag{CI/CD}
\cvtag{Confluence}
\cvtag{Jira}
\cvtag{Zenhub}
\cvtag{OpenSearch}
\cvtag{Sentry}
\cvtag{Linux RPM awareness}
\cvtag{BitBucket awareness}
\cvtag{Jenkins awareness}
\cvtag{Maven awareness}
\cvtag{Kubernetes awareness}
\cvtag{Réseaux awareness}
}

\divider\smallskip
\vspace{-.3cm}

\cvsubsection{Simulation, performance \& traitement du signal}

{\small
\cvtag{Traitement du signal}
\cvtag{Simulation numérique}
\cvtag{Analyse spectrale}
\cvtag{Analyse de bruit}
\cvtag{Signaux faibles}
\cvtag{Low-SNR}
\cvtag{Analyse de performance}
\cvtag{Performance système}
\cvtag{Modélisation}
\cvtag{Comparaison de modèles}
\cvtag{Analyse de sensibilité}
\cvtag{Data quality}
\cvtag{Validation de résultats}
\cvtag{Instrumentation}
\cvtag{Acoustique awareness}
\cvtag{Sonar awareness}
}

\divider\smallskip
\vspace{-.3cm}

\cvsubsection{Coordination logiciel/système}

{\small
\cvtag{Coordination technique}
\cvtag{Interfaces logiciel/système}
\cvtag{Interfaces équipes développement}
\cvtag{Interfaces équipes système}
\cvtag{Documentation technique}
\cvtag{Revues techniques}
\cvtag{Synthèse technique}
\cvtag{Support décisionnel}
\cvtag{Priorisation}
\cvtag{Backlog}
\cvtag{Communication technique}
\cvtag{Anglais technique}
\cvtag{Collaboration internationale}
}

\divider\smallskip
\vspace{-.3cm}

\cvsubsection{Domaines d’application}

{\small
\cvtag{Spatial}
\cvtag{Systèmes complexes}
\cvtag{LISA ESA/NASA}
\cvtag{CNES}
\cvtag{ESA}
\cvtag{Instrumentation scientifique}
\cvtag{Systèmes logiciels critiques}
\cvtag{Plateformes de simulation}
\cvtag{Services numériques}
\cvtag{Infrastructures connectées}
\cvtag{Défense awareness}
\cvtag{Systèmes sonar awareness}
}

\end{document}